\documentclass[11pt, oneside]{amsart}

\usepackage{geometry}
\geometry{letterpaper}

\usepackage{ned-common}

\begin{document}

\title{Bloom Filters}
\maketitle

\newcommand{\bigO}[0]{\mathcal{O}}

\section{Setup}

\begin{enumerate}

\item The goal is to do a fast, space-efficient set membership. Our
scenario is that we don't otherwise have enough main memory to store all
the items.

\item To save memory, you are willing to suffer some \emph{false
positives}: \texttt{set.contains?(key)} is allowed to return
\texttt{true} even if \texttt{key} was never stored in \texttt{set}.
However, you insist on no \emph{false negatives}:
\texttt{set.contains?(key)} must return \texttt{true} if \texttt{key}
was ever stored in \texttt{set}.

\item This setup especially makes sense if there is some further, but
more costly, processing that can eliminate the false positives. For
instance, if we want to check if a key is an database, we might first
check a probabilistic datastructure, and only on a match (maybe true
positive, maybe false positive) do we load a block of the database from
disk.

\item Another example comes from Akamai. They don't cache anything until
it is requested at least twice. So when a key is first queried, they put
it in a bloom filter. When it is queried a second time, only then do
they put it into their cache.

\end{enumerate}

\section{Hash Set}

\begin{enumerate}

\item A simple first approach is to use a hash set. You hash the key to
$k$ bits, and store the key in the corresponding bucket. With this
approach, you have neither any false negatives (which would not be
tolerated), nor do you have any false positives (because you always
check the key). However, you must store every key! That can take a lot
of memory.

\item You may reduce memory by storing only the hash for a key, but not
the key itself. This does not introduce false negatives, but it creates
the possibility for a false positive: when the hash of the query
collides with a hash of a previously stored (but different) key.

\item The false positive rate depends on the number of bits used for the
hash, and the number of items stored in the hash set. Let's call the
number of bits of hash $B$, and the number of hash values $m = 2^B$. The
probability that a pair of distinct keys will suffer a hash collision is
$\frac{1}{m} = 2^{-B}$.

\item What is the probability that a next key will hash collide with
\emph{any} of $n$ previously stored hashes? The probability can be
\emph{no worse} than $\epsilon \approx \frac{n}{m} = n 2^{-B}$. This is
a worst-case estimate, because it presumes that there were \emph{no}
hash collisions amongst the prior $n$ entries. To achieve a false
positive rate of $\epsilon$, we need no more than  $B \approx
\log_2{n/\epsilon}$ bits of hash per key. Total storage for $n$ items is
thus $\bigO(n \log n)$.

\item In fact, a nice fact about the hash-set approach is that we could
choose $B$ very conservatively: maybe $B = \log_2 N / \epsilon$, where
$N$ is the \emph{most} number of keys we could consider storing. Then
our hash set storage is $\bigO(n \log_2 N / \epsilon)$. This is linear
in the number of keys \emph{actually} stored, and logarithmic in the
most keys we would ever want to store. This could be helpful if the hash
set is online and is dynamically updated, and we don't know at the start
how many hashes we will actually want to store.

\item A downside to the hash set implementation is that it can have
somewhat high overhead. We need to have a resizable array of buckets,
which typically incurs some indirection/poor locality.

\end{enumerate}

\section{Bloom Filter Intuition}

\begin{enumerate}

\item How can we reduce the memory requirement for storing our hash set?
Well, notice that we are storing hash values in the hash set. Each has
takes $B \approx \log_2 n / \epsilon$ bits. But this hash value sort-of
duplicates information already implicit in the \emph{location} of the
hash value in the hash set. Maybe there is a way to eliminate this
``duplication.''

\item One idea is to avoid storing the low bits of the hash values,
since this is already implicit in the bucket location. You can just
store the top bits in the bucket. I believe this is more-or-less the
idea behind \define{quotient filters}. I kind of came up with this
myself. It should allow resizing. This works easiest bitwise if you
always use $2^k$ buckets, for some $k$ that can change.

\item Consider an alternative in which, at each bucket, we only store
one-bit: whether an item has been previously stored in the bucket. We
will pop a bit at index $i$ if an item with hash $i$ is stored. Because
we won't store hash values explicitly, we need a bit-array of $2^B$
positions. Of course, this uses \emph{even more} memory than our hash
set.

\item Still, this is already an example of a Bloom Filter. A Bloom
Filter will have two parameters: $m$, the number of bits in the
bit-array, and $k$, the number of hashes used to probe the bit-array. In
this example, we have $m = 2^B$, and $k=1$. We will stop worrying about
$B$ and focus on $m$, because we're just going to use a hash value to
uniformly select one of $m$ buckets to query or pop.

\item We will now explore whether we can tune the $m, k$ parameters to
get the same target false-positive rate but with less memory. Imagine
doubling $k$ from 1 to 2. Now, when adding to the bloom filter, you
calculate two hashes and you pop two bits. To query, you check two bits.

\item Let's consider what happens when you double the number of hash
functions. What happens to the probability of a false positive?

\begin{IEEEeqnarray*}{rClrCl}
  p_\text{bit}
\eqcol
  \frac{kn}{m}
&\quad
  p'_\text{bit}
\eqcol
  \frac{2kn}{m}
\\
  p_\text{k-query}
\eqcol
  \left( \frac{kn}{m} \right)^k
&\quad
  p'_\text{k-query}
\eqcol
  \left( \frac{2kn}{m} \right)^k
\end{IEEEeqnarray*}

\noindent
Then, because we probe double the positions:

\begin{nedqn}
  p'_\text{2k-query}
\eqcol
  \left( \frac{2kn}{m} \right)^{2k}
\\
\eqcol
  2^{2k}
  p_\text{bit}^{2k}
\end{nedqn}

\noindent
This suggests that you should continue to increase $k$ while
$p_\text{bit}$ is less than 50\%. So we believe that $k = \frac{m}{2n}$
is optimal.

\item In that case, what size $m$ should we pick to achieve a target
false positive rate of $\epsilon$? We want:

\begin{nedqn}
  p_\text{$k$-query}
\eqcol
  \left( \frac{kn}{m} \right)^k
= \epsilon
\\
  0.5 ^ {\frac{m}{2n}}
\eqcol
  \epsilon
\\
  \frac{m}{2n} \log_2{0.5}
\eqcol
  \log_2 \epsilon
\\
  m
\eqcol
  -2n \log_2 \epsilon
\intertext{thus we also know}
  k
\eqcol
  -\log_2 \epsilon
\end{nedqn}

\item This is \emph{basically} correct. However, it is too conservative
to believe that $kn/m$ bits will be popped on entry of $n$ items.
Probably there is \emph{some} overlap, which means that false-positive
rate of a query after $n$ insertions is actually lower than we plugged
in. When we correct for this, we get:

\begin{nedqn}
  m
\eqcol
  -n \frac{\ln \epsilon}{\ln(2)^2}
\\
  m
&\approx&
  -1.44 n \log_2 \epsilon
\intertext{and}
  k
\eqcol
  -\frac{\ln \epsilon}{\ln(2)}
\\
  k
\eqcol
  -\log_2 \epsilon
\end{nedqn}

\end{enumerate}

\section{Final Notes}

\begin{enumerate}

\item The false positive rate will be start at zero and grow (degrade)
as you add items. You can keep adding items until half the bits are
popped without pushing beyond the goal false positive rate. Indeed, you
can keep adding items even after reaching half of the bits popped. What
will happen is that the false positive rate will continue to degrade,
but of course the bloom filter will otherwise still function.

\item You can do union easily with Bloom filters. You can also do
intersection, with some error (since the same bit can be set in both
filters as a result of inserting different elements).

\item There is a simple formula to calculate the expected count in a
bloom filter. The naive version is $n = \frac{\text{num bits
popped}}{k}$, but this is too small, because of possibility of prior
collisions in the bit array.

\item There are attempts to allow you to delete, by keeping a count at
each array position. However, this is not easy and I think all attempts
to solve this have tradeoffs.

\end{enumerate}

\end{document}

% ## Cuckoo Filter

% Conceptually: add a fingerprint of the key into a hash set. This will
% allow probabilistic removal, btw. Chaining is done via cuckoo hashing,
% so lookup is O(1). One advantage is that you will only look in at most
% 2 locations for the fingerprint, whereas you need to look in `k`
% locations for a bloom filter.

% In a Cuckoo filter, it is common to have a bucket store multiple
% fingerprints.

% Sources:

% * https://www.cs.cmu.edu/~dga/papers/cuckoo-conext2014.pdf
% * https://news.ycombinator.com/item?id=13056726
